% GENERATED FILE. Edit canonical lessons, then run npm run generate:books.
% canonical-source-hash: fnv1a64:8141677ab5ebde89
% canonical-lessons: BN-C321-kingbodonti, BN-C321-rupkotha, BN-C321-doityo, BN-C321-jadu, BN-C321-ghumano

\chapter{Legends}
\label{ch:bn-legends}
\hlchaptermodality{321}

\begin{chapteropening}
\textbf{By the end of this chapter:} \emph{I can tell a legend or a fairy tale about giants and magic.}

Complete the last lesson of chapter 321: I can tell a legend or a fairy tale about giants and magic.
\end{chapteropening}

\section[\texorpdfstring{\bn{কিংবদন্তি} (kingbôdônti)}{কিংবদন্তি (kingbôdônti)}]{\bn{কিংবদন্তি} (kingbôdônti) — a legend}

\label{lesson:BN-C321-kingbodonti}



\begin{quote}
Before the new one: say the Bengali for a dome, then the Bengali for architecture.
\end{quote}

\subsection*{You'll want to know: \bn{কিংবদন্তি}}
\textbf{\bn{কিংবদন্তি}} — \emph{kingbôdônti} — \textquotedblleft{}a legend\textquotedblright{}.

\begin{quote}
\bn{গ্রামের} \bn{একটা} \bn{কিংবদন্তি} \bn{আছে} \bn{যে} \bn{পাহাড়ের} \bn{গুহায়} \bn{এক} \bn{রাজা} \bn{থাকতেন।}
\end{quote}

\begin{quote}
\emph{grāmer ekṭā kingbôdônti āchhe je pāhāṛer guhāy ek rājā thākten.}
\end{quote}

\begin{quote}
There's a village legend that a king once lived in the cave in the hills.
\end{quote}

\begin{cousinweb}
\textbf{\bn{কিংবদন্তি}} (\emph{kingbôdônti}) is Sanskrit \emph{kiṃvadantī}, 'what do they say?' — a common saying, a rumour, a tradition (Monier-Williams). A legend is what people say.
\end{cousinweb}

\begin{grammarlens}[title={What you've built}]
One more word for telling what happened.
\end{grammarlens}

\subsection*{Your turn}
\begin{itemize}
\raggedright
  \item \emph{Say it:} \emph{kingbôdônti}
  \item \emph{Say it:} \emph{kingbôdônti}, once more
  \item \emph{Say it:} \emph{grāmer ekṭā kingbôdônti āchhe je pāhāṛer guhāy ek rājā thākten.} — There's a village legend that a king once lived in the cave in the hills.
  \item \emph{Recall:} say the Bengali for a dome, then the Bengali for architecture, then say \emph{kingbôdônti} again
\end{itemize}

\begin{tcolorbox}[breakable,colback=teal!4,colframe=teal!35!black,title={Before you move on}]
What does \bn{কিংবদন্তি} mean? (\textquotedblleft{}A legend\textquotedblright{}.) Say it once more.
\end{tcolorbox}
\section[\texorpdfstring{\bn{রূপকথা} (rūpkôthā)}{রূপকথা (rūpkôthā)}]{\bn{রূপকথা} (rūpkôthā) — a fairy tale}

\label{lesson:BN-C321-rupkotha}



\begin{quote}
Before the new one: say the Bengali for architecture, then the Bengali for a legend.
\end{quote}

\subsection*{You'll want to know: \bn{রূপকথা}}
\textbf{\bn{রূপকথা}} — \emph{rūpkôthā} — \textquotedblleft{}a fairy tale\textquotedblright{}.

\begin{quote}
\bn{দিদিমা} \bn{রোজ} \bn{রাতে} \bn{আমাদের} \bn{একটা} \bn{রূপকথা} \bn{শোনাতেন।}
\end{quote}

\begin{quote}
\emph{didimā roj rāte āmāder ekṭā rūpkôthā shonāten.}
\end{quote}

\begin{quote}
Every night my grandmother used to tell us a fairy tale.
\end{quote}

\begin{cousinweb}
\textbf{\bn{রূপকথা}} (\emph{rūpkôthā}) is usually explained as \bn{রূপ}, 'form, beauty' — Sanskrit \emph{rūpa}, 'outward appearance' (Monier-Williams) — joined to \textbf{\bn{কথা}}, a story: a tale of wonders.
\end{cousinweb}

\begin{grammarlens}[title={What you've built}]
One more word for telling what happened.
\end{grammarlens}

\subsection*{Your turn}
\begin{itemize}
\raggedright
  \item \emph{Say it:} \emph{rūpkôthā}
  \item \emph{Say it:} \emph{rūpkôthā}, once more
  \item \emph{Say it:} \emph{didimā roj rāte āmāder ekṭā rūpkôthā shonāten.} — Every night my grandmother used to tell us a fairy tale.
  \item \emph{Recall:} say the Bengali for architecture, then the Bengali for a legend, then say \emph{rūpkôthā} again
\end{itemize}

\begin{tcolorbox}[breakable,colback=teal!4,colframe=teal!35!black,title={Before you move on}]
What does \bn{রূপকথা} mean? (\textquotedblleft{}A fairy tale\textquotedblright{}.) Say it once more.
\end{tcolorbox}
\section[\texorpdfstring{\bn{দৈত্য} (doityô)}{দৈত্য (doityô)}]{\bn{দৈত্য} (doityô) — a giant; a demon}

\label{lesson:BN-C321-doityo}



\begin{quote}
Before the new one: say the Bengali for a legend, then the Bengali for a fairy tale.
\end{quote}

\subsection*{You'll want to know: \bn{দৈত্য}}
\textbf{\bn{দৈত্য}} — \emph{doityô} — \textquotedblleft{}a giant; a demon\textquotedblright{}.

\begin{quote}
\bn{রূপকথায়} \bn{একটা} \bn{দৈত্য} \bn{রাজকুমারকে} \bn{বনে} \bn{আটকাল।}
\end{quote}

\begin{quote}
\emph{rūpkôthāy ekṭā doityô rājkumārke bône āṭkālô.}
\end{quote}

\begin{quote}
In the fairy tale a giant trapped the prince in the forest.
\end{quote}

\begin{cousinweb}
\textbf{\bn{দৈত্য}} (\emph{doityô}) is Sanskrit \emph{daitya}, a son of Diti, one of the demons of the old stories (Monier-Williams); in a Bengali fairy tale it is a huge, rough giant.
\end{cousinweb}

\begin{grammarlens}[title={What you've built}]
One more word for telling what happened.
\end{grammarlens}

\subsection*{Your turn}
\begin{itemize}
\raggedright
  \item \emph{Say it:} \emph{doityô}
  \item \emph{Say it:} \emph{doityô}, once more
  \item \emph{Say it:} \emph{rūpkôthāy ekṭā doityô rājkumārke bône āṭkālô.} — In the fairy tale a giant trapped the prince in the forest.
  \item \emph{Recall:} say the Bengali for a legend, then the Bengali for a fairy tale, then say \emph{doityô} again
\end{itemize}

\begin{tcolorbox}[breakable,colback=teal!4,colframe=teal!35!black,title={Before you move on}]
What does \bn{দৈত্য} mean? (\textquotedblleft{}A giant; a demon\textquotedblright{}.) Say it once more.
\end{tcolorbox}
\section[\texorpdfstring{\bn{জাদু} (jādu)}{জাদু (jādu)}]{\bn{জাদু} (jādu) — magic}

\label{lesson:BN-C321-jadu}



\begin{quote}
Before the new one: say the Bengali for a fairy tale, then the Bengali for a giant.
\end{quote}

\subsection*{You'll want to know: \bn{জাদু}}
\textbf{\bn{জাদু}} — \emph{jādu} — \textquotedblleft{}magic\textquotedblright{}.

\begin{quote}
\bn{জাদু} \bn{দিয়ে} \bn{দৈত্যটা} \bn{রাজকুমারকে} \bn{একটা} \bn{পাখি} \bn{বানাল।}
\end{quote}

\begin{quote}
\emph{jādu diye doityôṭā rājkumārke ekṭā pākhi bānālô.}
\end{quote}

\begin{quote}
With magic the giant turned the prince into a bird.
\end{quote}

\begin{cousinweb}
\textbf{\bn{জাদু}} (\emph{jādu}) is usually traced to Persian \emph{jādū}, 'magic, sorcery'. \bn{জাদুকর} is a magician.
\end{cousinweb}

\begin{grammarlens}[title={What you've built}]
One more word for telling what happened.
\end{grammarlens}

\subsection*{Your turn}
\begin{itemize}
\raggedright
  \item \emph{Say it:} \emph{jādu}
  \item \emph{Say it:} \emph{jādu}, once more
  \item \emph{Say it:} \emph{jādu diye doityôṭā rājkumārke ekṭā pākhi bānālô.} — With magic the giant turned the prince into a bird.
  \item \emph{Recall:} say the Bengali for a fairy tale, then the Bengali for a giant, then say \emph{jādu} again
\end{itemize}

\begin{tcolorbox}[breakable,colback=teal!4,colframe=teal!35!black,title={Before you move on}]
What does \bn{জাদু} mean? (\textquotedblleft{}Magic\textquotedblright{}.) Say it once more.
\end{tcolorbox}
\section[\texorpdfstring{\bn{ঘুমানো} (ghumāno)}{ঘুমানো (ghumāno)}]{\bn{ঘুমানো} (ghumāno) — to sleep}

\label{lesson:BN-C321-ghumano}



\begin{quote}
Before the new one: say the Bengali for a giant, then the Bengali for magic.
\end{quote}

\subsection*{You'll want to know: \bn{ঘুমানো}}
\textbf{\bn{ঘুমানো}} — \emph{ghumāno} — \textquotedblleft{}to sleep\textquotedblright{}.

\begin{quote}
\bn{জাদুর} \bn{জন্য} \bn{রাজকুমার} \bn{একশো} \bn{বছর} \bn{ঘুমালেন।}
\end{quote}

\begin{quote}
\emph{jādur jônyô rājkumār eksho bôchhôr ghumālen.}
\end{quote}

\begin{quote}
Because of the magic the prince slept for a hundred years.
\end{quote}

Past (he/she …): \bn{সেই} \bn{শিশু} \bn{ঘুমাল}, \emph{sei shishu ghumālô}.

\begin{cousinweb}
\textbf{\bn{ঘুমানো}} (\emph{ghumāno}) is filed in CDIAL 4485 under a Middle Indo-Aryan \emph{ghummati}, 'revolves', which lists Bengali \emph{ghumā}, 'to doze'; how the sense moved is not set out there.
\end{cousinweb}

\begin{grammarlens}[title={What you've built}]
One more word for telling what happened.
\end{grammarlens}

\subsection*{Your turn}
\begin{itemize}
\raggedright
  \item \emph{Say it:} \emph{ghumāno}
  \item \emph{Say it:} \emph{ghumāno}, once more
  \item \emph{Say it:} \emph{jādur jônyô rājkumār eksho bôchhôr ghumālen.} — Because of the magic the prince slept for a hundred years.
  \item \emph{Recall:} say the Bengali for a giant, then the Bengali for magic, then say \emph{ghumāno} again
\end{itemize}

\begin{tcolorbox}[breakable,colback=teal!4,colframe=teal!35!black,title={Before you move on}]
What does \bn{ঘুমানো} mean? (\textquotedblleft{}To sleep\textquotedblright{}.) Say it once more.
\end{tcolorbox}
